\documentclass[11pt]{article}
\usepackage[margin=1in]{geometry}
\usepackage{amsmath,amssymb,amsthm}
\usepackage{graphicx}
\usepackage[colorlinks=true,linkcolor=blue,citecolor=blue,urlcolor=blue]{hyperref}

\newtheorem{theorem}{Theorem}
\newtheorem{lemma}{Lemma}
\newtheorem{corollary}{Corollary}
\newtheorem{remark}{Remark}

\newcommand{\R}{\mathbb R}
\newcommand{\C}{\mathbb C}
\newcommand{\N}{\mathbb N}
\newcommand{\dens}{\operatorname{dens}}
\newcommand{\cf}{\operatorname{cf}}
\newcommand{\card}{\operatorname{card}}
\newcommand{\spanop}{\operatorname{span}}

\title{An Explicit ZFC Cardinal Obstruction to Laugwitz's Complete-Norm Question}
\author{Automatic Math Discovery Smoke Test}
\date{July 1, 2026}

\begin{document}
\maketitle

\section*{Source and Extracted Question}

The source is Renan J. S. Isneri, Josias V. Baca, and Lucas M. Fernandes,
\emph{Norms and Non-Equivalence in Infinite-Dimensional Banach Spaces},
arXiv:2508.15686.  In the introduction the authors ask:

\begin{quote}
Does there exist some cardinal \(\kappa>\aleph_0\) for which no vector space
with a Hamel basis of that cardinality can be endowed with a norm that turns it
into a Banach space?
\end{quote}

They attribute the question to Laugwitz \cite{Laugwitz} and state that no
explicit example of such a cardinal is currently known.

\begin{center}
\includegraphics[width=\textwidth]{figures/open_problem_crop.png}
\end{center}

\section*{Answer}

The cardinal
\[
  \kappa=\aleph_\omega
\]
is an explicit ZFC example.  No real or complex vector space with Hamel basis
of cardinality \(\aleph_\omega\) admits a norm whose metric completion is
already the given space.

\section*{Idea of the Proof}

There are two independent obstructions to a cardinal being a Banach-space
Hamel dimension.

First, a theorem of Halbeisen and Hungerbuehler, cited by the source paper as
\cite[Theorem 3.5]{HalbeisenHungerbuehler}, says that every Hamel basis of an
infinite-dimensional real or complex Banach space \(X\) has cardinality
\(|X|\).  Thus a Banach Hamel dimension cannot be below the continuum.

Second, if \(X\) has density character \(\delta\), then the metric structure
forces
\[
  |X|=\delta^{\aleph_0}.
\]
Consequently any infinite-dimensional Banach Hamel dimension \(\kappa\) must
satisfy \(\kappa^{\aleph_0}=\kappa\).  But by Koenig's theorem,
\[
  \aleph_\omega^{\aleph_0}>\aleph_\omega.
\]
This excludes \(\aleph_\omega\) whenever the continuum is at most
\(\aleph_\omega\), and the first obstruction excludes it whenever the
continuum is larger.

\section*{Cardinality Lemma for Banach Spaces}

\begin{lemma}
\label{lem:banach-cardinality}
Let \(X\) be an infinite-dimensional Banach space over \(\R\) or \(\C\), and
let \(\delta=\dens(X)\).  Then
\[
  |X|=\delta^{\aleph_0}.
\]
\end{lemma}

\begin{proof}
Let \(D\subset X\) be dense with \(|D|=\delta\).  Every \(x\in X\) is the
limit of some sequence from \(D\), so
\[
  |X|\leq |D|^{\aleph_0}=\delta^{\aleph_0}.
\]

For the reverse inequality we first produce a separated family in the unit
sphere of size \(\delta\).  By transfinite induction choose unit vectors
\((x_\alpha)_{\alpha<\delta}\) so that
\[
  \operatorname{dist}\bigl(x_\alpha,
    \overline{\spanop}\{x_\beta:\beta<\alpha\}\bigr)>1/2.
\]
This is possible because the closed linear span of fewer than \(\delta\)
vectors has density \(<\delta\), hence is not all of \(X\); Riesz's lemma then
supplies the next vector.  Thus
\[
  \|x_\alpha-x_\beta\|>1/2\qquad(\alpha\neq\beta).
\]

For each sequence \(s=(\alpha_n)_{n\geq 1}\in\delta^\N\), define
\[
  \Phi(s)=\sum_{n=1}^{\infty}6^{-n}x_{\alpha_n}\in X.
\]
The series converges absolutely.  If \(s\neq t\), let \(n\) be the first index
with \(\alpha_n\neq\beta_n\).  Then
\[
\begin{aligned}
  \|\Phi(s)-\Phi(t)\|
  &\geq 6^{-n}\|x_{\alpha_n}-x_{\beta_n}\|
      -2\sum_{m>n}6^{-m}  \\
  &> \frac12\,6^{-n}-\frac25\,6^{-n}>0.
\end{aligned}
\]
So \(\Phi\) is injective and \(|X|\geq\delta^{\aleph_0}\).
\end{proof}

\section*{The Obstruction}

\begin{theorem}
\label{thm:aleph-omega-obstruction}
Let \(E\) be a real or complex vector space with Hamel dimension
\(\aleph_\omega\).  No norm on \(E\) makes \(E\) into a Banach space.
\end{theorem}

\begin{proof}
Suppose that some norm makes \(E\) a Banach space \(X\).

Let \(\mathfrak c=2^{\aleph_0}\).  If \(\mathfrak c>\aleph_\omega\), then the
underlying set of a real or complex vector space with Hamel dimension
\(\aleph_\omega\) has cardinality \(\mathfrak c\).  But
Halbeisen--Hungerbuehler's theorem gives
\[
  \dim_{\mathrm H}(X)=|X|=\mathfrak c,
\]
contradicting \(\dim_{\mathrm H}(X)=\aleph_\omega\).

It remains to consider the case \(\mathfrak c\leq\aleph_\omega\).  Then
\(|X|=\aleph_\omega\), and again
Halbeisen--Hungerbuehler gives
\[
  \aleph_\omega=\dim_{\mathrm H}(X)=|X|.
\]
By Lemma \ref{lem:banach-cardinality}, for
\(\delta=\dens(X)\) one has
\[
  \aleph_\omega=|X|=\delta^{\aleph_0}.
\]
Therefore
\[
  \aleph_\omega^{\aleph_0}
  =(\delta^{\aleph_0})^{\aleph_0}
  =\delta^{\aleph_0}
  =\aleph_\omega.
\]
This contradicts Koenig's theorem, since \(\cf(\aleph_\omega)=\aleph_0\) and
therefore
\[
  \aleph_\omega^{\aleph_0}>\aleph_\omega.
\]
The contradiction proves the claim.
\end{proof}

\begin{corollary}
Every singular cardinal \(\kappa\) of countable cofinality gives the same
kind of obstruction, after the identical case split for the continuum:
if \(\kappa\geq\mathfrak c\), Koenig's theorem gives
\(\kappa^{\aleph_0}>\kappa\), and if \(\kappa<\mathfrak c\), the
Halbeisen--Hungerbuehler cardinality equality already rules it out.
\end{corollary}

\section*{Verification and Novelty Notes}

The proof is set-theoretic and does not require numerical or symbolic
computation.  The local run indexes were searched for the arXiv id
\texttt{2508.15686}, the phrases \texttt{Laugwitz},
\texttt{Hamel bases of Banach spaces}, and \texttt{aleph\_omega}, and no
duplicate solution packet was found.  A local source search found the source
paper's citation of Halbeisen--Hungerbuehler and an independent Banach-space
paper using the standard consequence
\(\aleph_\omega^{\aleph_0}>\aleph_\omega\) of Koenig's theorem, but no packet
answering the present Laugwitz question.

Bounded web searches for the arXiv id, the title, and close phrases involving
Laugwitz, Hamel bases, Banach spaces, and \(\aleph_\omega\) found no later
explicit answer.  Because the proof uses classical ingredients, the
mathematical claim should be checked by a human reviewer and may be
reclassified as a literature-implied answer if an exact prior formulation is
located.

\section*{References}

\begin{thebibliography}{9}

\bibitem{ArendtNittka}
W. Arendt and R. Nittka,
\emph{Equivalent complete norms and positivity},
Archiv der Mathematik \textbf{92} (2009), no. 5, 414--427.

\bibitem{HalbeisenHungerbuehler}
L. Halbeisen and N. Hungerbuehler,
\emph{The cardinality of Hamel bases of Banach spaces},
East-West Journal of Mathematics \textbf{2} (2000), no. 2, 153--159.

\bibitem{Laugwitz}
D. Laugwitz,
\emph{\"Uber vollst\"andige Normtopologien in linearen R\"aumen},
Archiv der Mathematik \textbf{6} (1955), no. 2, 128--131.

\bibitem{Source}
R. J. S. Isneri, J. V. Baca, and L. M. Fernandes,
\emph{Norms and Non-Equivalence in Infinite-Dimensional Banach Spaces},
arXiv:2508.15686.

\end{thebibliography}

\end{document}
